\documentclass{situationist-preprint}

\title{The Fittingness of Resentment}
\msid{ACT-09-a}
\msfamily{ACT-09}
\msdiscipline{Action, Agency and Free Will}
\mstype{Principal manuscript}
\mskeywords{resentment; reactive attitudes; fittingness; anger; quality of will; standing to blame; Strawson}

\begin{document}
\maketitle

\begin{abstract}
Resentment has had bad press. Philosophers from Seneca to Nussbaum have held that it is irrational, since it wishes harm on the wrongdoer and the harm repairs nothing, and skeptics about free will hold that it presupposes a desert that no one has. We argue that resentment is justified in the sense that it is a fitting response to wrongdoing: when an agent wrongs us with ill will or culpable disregard, resentment correctly represents what she has done. The argument rests on a view of emotions as evaluative representations and on Strawson's account of the reactive attitudes as responses to the quality of another's will. Fittingness is a matter of correctness, not of usefulness, and it does not depend on whether feeling or showing resentment does any good. We then separate the fittingness of the emotion from the norms that govern its expression. Expressing resentment is an action, subject to considerations of proportion, standing and consequence, and these may make its expression wrong when the feeling is fitting. The separation answers the main objections to resentment while preserving what they get right.
\end{abstract}

\section{The problem}

\citet{Strawson1962} argued that resentment, indignation and gratitude are \emph{reactive attitudes}: responses to the quality of will that other people show towards us and towards others, whose presence is constitutive of ordinary interpersonal relationships. The question whether these attitudes are justified, he held, cannot be settled by a theory of determinism, since the attitudes answer to whether another has shown us good will, ill will or indifference, and not to whether her will was determined. Strawson's argument was directed at the threat from determinism. But resentment faces older and more direct objections. \citet{Butler1726}, in his sermon on resentment, defended it as a natural guard against injury, while warning against its abuse. The Stoic tradition held it irrational. \citet{Nussbaum2016} argued that anger in the standard case involves a wish that the wrongdoer suffer, which is either irrational, since the suffering repairs nothing, or a concern with relative status, which is ignoble; the rational response to wrongdoing is a transition to concern for future welfare. \citet{Pereboom2001} argued that resentment presupposes that the wrongdoer basically deserves to suffer, which requires a freedom no one has. Defenders of resentment have replied that it is a form of protest \citep{Hieronymi2001}, that it embodies respect for oneself and for the wrongdoer as a responsible agent \citep{MurphyHampton1988}, and that it can be apt even when expressing it is counterproductive \citep{Srinivasan2018}.

We argue that resentment is a fitting response to wrongdoing, and that its expression is subject to further norms. The argument relies on a distinction, drawn by \citet{DArmsJacobson2000}, between the fittingness of an emotion, which is a matter of whether its evaluative presentation of its object is correct, and other kinds of appropriateness, moral or prudential.

\section{Framework}

\begin{definition}[Resentment]\label{d:resent}
\emph{Resentment} is the emotion in which a subject is aggrieved by an action of another agent, presenting it as a wrong done to her that expresses the agent's ill will or culpable disregard for her standing, together with a disposition to protest the wrong and to demand acknowledgment of it.
\end{definition}

The definition follows \citet{Strawson1962} in making resentment a response to the quality of the other's will, and \citet{Wallace1994} in tying it to the violation of expectations that we hold one another to. Indignation is the corresponding attitude felt on another's behalf. The disposition to protest and to demand acknowledgment is the motivational profile of resentment; a wish that the wrongdoer suffer may accompany it, but is not part of it.

\begin{principle}[Evaluative content]\label{p:content}
Emotions represent evaluative properties of their objects. An emotion is \emph{fitting} if and only if its object has the evaluative property it represents it as having.
\end{principle}

\begin{principle}[Kinds of reason]\label{p:kinds}
That an emotion is fitting is a reason of the right kind for having it. That having or expressing it would have good or bad effects is a reason of the wrong kind with respect to its fittingness: it bears on whether it is good to have or express the emotion, not on whether the emotion is correct.
\end{principle}

\begin{assumption}[Responsible agency]\label{a:resp}
Some agents are responsible for some of their actions, in the sense that their actions issue from mechanisms that are their own and reasons-responsive, and through which their quality of will is expressed.
\end{assumption}

Principle~\ref{p:content} is the thesis that MND-17 defends: fear represents danger, and emotions in general have evaluative content. It gives a clear sense to the question whether resentment is justified: it is justified, in the first instance, if it is fitting, that is, if what it represents is so. Principle~\ref{p:kinds} is the distinction of \citet{DArmsJacobson2000}: envy of a rival's success may be fitting although it is ignoble to feel it, and fear of a harmless spider unfitting although it is useful in preparing a child for dangerous ones. Assumption~\ref{a:resp} is the conclusion of ACT-01, which argues that free will is real and compatible with determinism and indeterminism alike.

\section{The argument}

\begin{lemma}\label{l:represents}
Resentment of an action $a$ by agent $S$ represents $a$ as (i) a wrong, (ii) done to the resenting subject, (iii) by a responsible agent, (iv) expressing her ill will or culpable disregard.
\end{lemma}

\begin{proof}
By Definition~\ref{d:resent} and Principle~\ref{p:content}, resentment presents its object as a wrong done to the subject that expresses ill will or culpable disregard; these are clauses (i), (ii) and (iv). Ill will and culpable disregard are qualities of will that only a responsible agent can express: an infant or a person in the grip of psychosis may harm us, but cannot disregard us culpably. This is the observation behind Strawson's account of the excusing and exempting conditions, which inhibit resentment by showing that the agent did not express the quality of will it presupposes \citep{Strawson1962}. So resentment also represents $S$ as a responsible agent, clause (iii).
\end{proof}

\begin{theorem}\label{t:main}
When a responsible agent wrongs someone, through ill will or culpable disregard, the victim's resentment of the wrong is fitting.
\end{theorem}

\begin{proof}
Suppose a responsible agent $S$ wrongs $V$ through ill will or culpable disregard. By Lemma~\ref{l:represents}, $V$'s resentment represents $S$'s action as a wrong, done to $V$, by a responsible agent, expressing ill will or culpable disregard. Each of these is true by supposition. By Principle~\ref{p:content}, the resentment is fitting. Assumption~\ref{a:resp} ensures that the supposition is satisfiable: there are responsible agents whose wrongs express their quality of will, so the theorem is not vacuous.
\end{proof}

\begin{corollary}\label{c:indig}
The same holds for indignation: when a responsible agent wrongs a third party through ill will or culpable disregard, indignation on the third party's behalf is fitting.
\end{corollary}

\begin{proof}
Indignation represents what resentment represents, with clause (ii) replaced by the clause that the wrong is done to someone. The proof of Theorem~\ref{t:main} applies with this change.
\end{proof}

The theorem gives the sense in which resentment is justified that is relevant to the question whether there is anything wrong with it as such. It is not justified in every case: resentment of an excused harm, of an imagined slight, or of an agent who did not act from ill will is unfitting, because it misrepresents. And fitting resentment may be disproportionate in intensity, when it represents a minor wrong as grave. But in the central case, the response of a victim to a culpable wrong, it represents things as they are.

\section{The norms of expression}\label{s:norms}

\begin{proposition}\label{p:express}
That resentment is fitting does not settle whether it is permissible or advisable to express it, or to express it in a given way.
\end{proposition}

\begin{proof}
Expressing resentment is an action: a reproach, a demand for apology, a withdrawal of goodwill, a public complaint. Whether an action is permissible or advisable depends on the practical reasons that bear on it, including its effects on the wrongdoer, the victim and others, and on the relationships in which the parties stand. By Principle~\ref{p:kinds}, these reasons do not bear on the fittingness of the emotion, and the fittingness of the emotion is not among the considerations that settle by itself what one ought to do. So the fittingness of resentment leaves the question of its expression open.
\end{proof}

Three further norms govern expression. The first is proportion: the expression should not represent the wrong as graver than it was, nor impose burdens on the wrongdoer out of proportion to it. The second is standing: a person who has done the same wrong without acknowledging it lacks standing to reproach another for it, though her resentment may still be fitting. The third is consequence: when expression would provoke escalation, endanger the victim or harm innocent parties, there may be decisive reason not to express it, or to express it differently. These norms explain why the victim of a slight may rightly keep her resentment to herself, and why the expression of fitting resentment can itself be a wrong. They also explain the role of forgiveness, which the companion manuscript ACT-09-b takes up: when a wrongdoer has acknowledged the wrong and repudiated it, the wrong no longer expresses her present will toward the victim in the same way, and the reasons for continuing to protest it lapse.

The account fits the use ETH-26-a makes of resentment in its theory of punishment as communication. If resentment of a wrong done to oneself is fitting, an offender who would fittingly resent the same wrong done to her cannot consistently deny that censure of her own wrong is warranted. The censure that punishment communicates is, on that account, an expression of the indignation the community fittingly feels, and it is the state's expression, not the feeling, that is subject to the proportionality and termination conditions ETH-26 derives.

\section{Objections and replies}

\begin{objection}[Payback]
Resentment includes a wish that the wrongdoer suffer. Suffering does not undo the wrong. So resentment is irrational in its aim, whatever it represents \citep{Nussbaum2016}.
\end{objection}

\begin{reply}
The objection concerns the motivational profile of resentment, not its representational content, and on Definition~\ref{d:resent} the characteristic aim of resentment is protest and acknowledgment, not suffering. \citet{Hieronymi2001} argues that resentment protests a claim implicit in the wrong, that the victim may be treated so, and that this claim persists as long as the wrong stands unacknowledged. The aim of such protest is not achieved by the wrongdoer's suffering but by her acknowledgment. Where resentment includes a wish for the wrongdoer's suffering beyond this, that wish is subject to the norms of Section~\ref{s:norms}, and may be one that the resentful person should not act on. Its presence in some cases does not make the emotion unfitting.
\end{reply}

\begin{objection}[Basic desert]
Resentment presupposes that the wrongdoer deserves to suffer in a basic sense, and basic desert requires a freedom that is incompatible with determinism and with indeterminism alike \citep{Pereboom2001}.
\end{objection}

\begin{reply}
By Lemma~\ref{l:represents}, resentment presupposes that its target is a responsible agent who expressed ill will or culpable disregard. It does not represent the target as deserving suffering, since its content concerns the wrong and the will behind it, not a response to be inflicted. The presupposition of responsible agency is satisfied, on the argument of ACT-01-a, by agents whose actions issue from their own reasons-responsive mechanisms, and that argument shows such agency compatible with determinism and with indeterminism. ACT-08 adds that responsibility does not require that the agent be the ultimate source of her action.
\end{reply}

\begin{objection}[Uselessness]
Resentment makes the resentful person miserable and her relationships worse; a rational person would do without it.
\end{objection}

\begin{reply}
These are reasons of the wrong kind (Principle~\ref{p:kinds}). They may give a person reason to cultivate a disposition to feel less resentment, or to express it less, but they do not show that resentment misrepresents its object. \citet{Srinivasan2018} has pressed the point in the case of victims of injustice who are told that their anger is counterproductive: the anger may be counterproductive and still apt, and to demand that the victim choose between fitting feeling and prudent action is itself a burden. Moreover, the consideration is not one-sided: \citet{Butler1726} and \citet{MurphyHampton1988} argue that resentment protects self-respect, and a person incapable of it fails to register wrongs done to her as wrongs.
\end{reply}

\begin{objection}[Hypocrisy]
Most of us have done the wrongs we resent in others. Resentment from such persons is unjustified.
\end{objection}

\begin{reply}
The hypocrite lacks standing to reproach, and so has reason not to express her resentment, or to express it together with an acknowledgment of her own fault. But her resentment represents the other's wrong correctly; her own wrong does not make his any less a wrong done to her. Standing is a norm of expression, and the objection confirms the separation of Proposition~\ref{p:express}.
\end{reply}

\section{Conclusion}

Resentment represents another's action as a wrong done to the subject by a responsible agent through ill will or culpable disregard. When things are as it represents them, it is fitting, and its fittingness does not depend on its usefulness or on a desert that outruns responsibility. Its expression is an action, subject to norms of proportion, standing and consequence that may tell against expressing a fitting resentment. Resentment is justified as a response to wrongdoing; what one does with it is a further question.

\begin{thebibliography}{99}\small
\bibitem[Butler(1726)]{Butler1726} Butler, J. (1726). \emph{Fifteen Sermons Preached at the Rolls Chapel}. Sermon VIII, Upon resentment. London.
\bibitem[D'Arms and Jacobson(2000)]{DArmsJacobson2000} D'Arms, J. and Jacobson, D. (2000). The moralistic fallacy: on the \enquote{appropriateness} of emotions. \emph{Philosophy and Phenomenological Research}, 61(1), 65--90.
\bibitem[Hieronymi(2001)]{Hieronymi2001} Hieronymi, P. (2001). Articulating an uncompromising forgiveness. \emph{Philosophy and Phenomenological Research}, 62(3), 529--555.
\bibitem[Murphy and Hampton(1988)]{MurphyHampton1988} Murphy, J.~G. and Hampton, J. (1988). \emph{Forgiveness and Mercy}. Cambridge University Press.
\bibitem[Nussbaum(2016)]{Nussbaum2016} Nussbaum, M.~C. (2016). \emph{Anger and Forgiveness: Resentment, Generosity, Justice}. Oxford University Press.
\bibitem[Pereboom(2001)]{Pereboom2001} Pereboom, D. (2001). \emph{Living Without Free Will}. Cambridge University Press.
\bibitem[Srinivasan(2018)]{Srinivasan2018} Srinivasan, A. (2018). The aptness of anger. \emph{The Journal of Political Philosophy}, 26(2), 123--144.
\bibitem[Strawson(1962)]{Strawson1962} Strawson, P.~F. (1962). Freedom and resentment. \emph{Proceedings of the British Academy}, 48, 1--25.
\bibitem[Wallace(1994)]{Wallace1994} Wallace, R.~J. (1994). \emph{Responsibility and the Moral Sentiments}. Harvard University Press.
\end{thebibliography}

\end{document}
